\section{Propagation through the original inverse fibres and spatial flow}
\label{sec:x4-spatial}

Use the exact per-channel arithmetic heat function $\mathcal K_t(s)$
defined by \eqref{xat:K}, with all source constants in
\eqref{x4:distribution}. Retain the incidence count by setting
$G_t(s)=a_{ef}\mathcal K_t(s)$. Define
\[
g_t(r)=G_t(-r^2/2),\qquad
\Psi^{(4)}(t,q)=G_t(\sigma(q)),\quad\sigma(q)=F_1(q)/2.
\]
These are entire compositions of the actual transformed arithmetic
function; they retain the incidence factor $a_{ef}$. Since the new
transform commutes with the actual Newman heat evolution, the already
proved coordinate chain rules apply to it verbatim and give
\begin{align}
\partial_tg_t&=
 -\frac1{4r^2}\partial_r^2g_t+\frac1{4r^3}\partial_rg_t
                       &&(r\ne0),\notag\\
\partial_t\Psi^{(4)}+\tfrac14U^2\Psi^{(4)}&=0,\notag\\
\Delta_q\Psi^{(4)}&=N(q)G_{ss}(t,\sigma(q))+
                              L(q)G_s(t,\sigma(q)).
\label{x4:spatial-equations}
\end{align}
Every coefficient $N,L,U$ is the original one proved in
Section~\ref{sec:eo}; no Euclidean diffusion coefficient has been
replaced by the root-coordinate generator. For any retained time map
$t=\theta(\tau)$, the full residual is explicitly
\[
(\partial_\tau+U\cdot\nabla_q-\nu\Delta_q)
 G_{\theta(\tau)}(\sigma(q))
=\left(-\frac{\theta'}4-\nu N\right)G_{ss}
                         +(1-\nu L)G_s,
\]
with the derivatives on the right evaluated at
$(\theta(\tau),\sigma(q))$. The proof is the time chain rule and the
three identities in \eqref{x4:spatial-equations}, so the force-sized
remainder is calculated rather than assumed absent.

The three spectral translations also have exact maps in the retained
original root coordinates. For
$\lambda\in\{3,9/2,13/2\}$ define the entire algebraic correspondence
\[
\mathcal C_\lambda=\{(r,u)\in\C^2:u^2=r^2-2\lambda\}.
\]
It satisfies $-u^2/2=-r^2/2+\lambda$. Consequently
\[
Z_t(-r^2/2+\lambda)=f_t(u),\qquad
Z_s(t,-r^2/2+3)=-\frac{\partial_uf_t(u)}u
                 \quad ((r,u)\in\mathcal C_3,\ u\ne0),
\]
where $f_t(u)=Z_t(-u^2/2)$ is the original function, already proved
even and entire. The second formula follows from
$\partial_uf_t=-uZ_s(t,-u^2/2)$; thus its quotient has a unique
entire continuation at $u=0$ equal to $-Z_s(t,0)$. Neither formula
depends on the choice of sign of $u$.

In these actual correspondences the full propagated function is
\[
g_t(r)=a_{ef}\left[c_0f_t(u_3)+
 \beta\frac{\partial_uf_t(u_3)}{u_3}
 +c_1f_t(u_{9/2})+c_2f_t(u_{13/2})\right],\qquad
u_\lambda^2=r^2-2\lambda.
\]
Each potentially branched expression descends to the displayed entire
function of $r$, by the proven evenness and removable quotient. The
finite inverse point for a nonzero $u_\lambda$ is exactly
$\Gamma(u_\lambda)=(-1/(2u_\lambda),3u_\lambda,26u_\lambda^2)$;
its target is
\[
F(\Gamma(u_\lambda))=(-r^2+2\lambda,0,0).
\]
At $u_\lambda=0$ the completed root point still has no finite inverse
in the $x\ne0$ chart, as its unchanged first coordinate shows. The
additional finite point with $x=0$ is retained by the earlier complete
fibre calculation. The three new branch values are precisely
$r^2=6,9,13$; they arise from the source's three specified energies.

All of this also retains the full original cubic coefficient family.
For fixed $b,c$ and $s(r)=(cr^3-2r^2+br)/4$, set
$g_{t;b,c}(r)=G_t(s(r))$. The exact same chain rule gives
\[
\partial_tg_{t;b,c}=
 -\frac4{P_r^2}\partial_r^2g_{t;b,c}
 +\frac{4P_{rr}}{P_r^3}\partial_rg_{t;b,c}
 \quad(P_r=3cr^2-4r+b\ne0).
\]
Thus the source fourth coefficient is propagated into the coefficient
family as well as its $b=c=0$ slice. This is a completed map of the
source observable and its actual heat dynamics. For the per-channel
function, recovery of the original $\Xi$ uses the full-domain inverse
automorphism proved in Section~\ref{sec:x4-arithmetic}. For any actual
pair with $a_{ef}>0$, the recovered input is
$Z_t=\mathcal T_{\mathcal D}^{-1}(G_t/a_{ef})$; the explicit separated
pair in the source proof has $a_{ef}\ge1$. At $a_{ef}=0$ this fourth
coefficient is exactly zero and carries no pair observation. A zero of a finite linear
combination is not silently substituted for a zero of its recovered
arithmetic input.
